\clearpage
\fancyhead[L]{\small GAUSS TO VON LINDENAU.}
\hypertarget{native-316-lower}{}
\begin{center}
\textbf{GAUSS TO VON LINDENAU.}

\par\vspace{4pt}\hrule\vspace{4pt}
\fontsize{10}{12}\selectfont Monatliche Correspondenz. Vol. XIX. Pp. 501..503. 1809. May.
\par\vspace{4pt}\hrule\vspace{4pt}
\end{center}

\begin{flushright}Göttingen, 20 May 1809.\end{flushright}

\par\indent I enclose the results of my calculations for improving the elements of the new planets. I completed some of these calculations a while ago. For Ceres, you will notice some differences between my results and \textsc{Santini}’s (Monatl. Corresp. Feb. p. 191). This is probably due in part to \textsc{Santini}’s neglecting the perturbations when he compared the \emph{individual} observations with the XI elements. I cannot approve of that, because these perturbations by no means have a uniform effect on the geocentric positions over the 22 days, so the mean difference must not be taken as the true one. Mr. \textsc{Santini} also seems to have made some errors in his subsequent calculation of the perturbations. From my tables I find a perturbation in longitude of $+6'\;6''4$ and one in the radius vector of $+0{,}001570$. You will notice that these figures bear almost no resemblance to \textsc{Santini}’s. \SourceLeaf{317}The perturbation in latitude, which of course cannot be compared with \textsc{Santini}’s because I used more complete equations, is $+56''8$. — On the enclosed sheet I also send you the ephemerides of Ceres, Juno, and Vesta for their next periods of visibility. Professor \textsc{Harding} calculated the one for Juno from the elements previously printed in the Monatl. Corresp. The ephemeris for Vesta, based on the IV elements, is by Dr. \textsc{Schumacher}, who is now staying with us. I undertook the one for Ceres myself. [\textsc{Von Lindenau} included this ephemeris in the article “Fortgesetzte Nachrichten u.\,s.\,f.” in the May 1809 issue of Mon. Corr., pages 504 to 515. It gives geocentric right ascension and declination to the minute and the logarithm of the distance to four decimal places, at four-day intervals for midnight in Göttingen from 18 July 1809 to 16 March 1810.] Only Pallas is still missing. Unfortunately I have yet to see any usable observations of it from last year, so I have been unable to make any calculations for it.

\Needspace{23\baselineskip}
\par\indent Regarding the notice of my programme in the February issue [in which an unnamed author “briefly presented in this journal the main formulas on which the method rests” and kept the numbering of the formulas and the notation from \textsc{Gauss}’s Methodus peculiaris elevationem poli determinandi, page 39 and the following pages of this volume, unchanged], I should also point out that it wrongly speaks of equal altitudes (p. 136). The problem has no such condition at all. The following addition is also worth making. “When $\lambda$ is determined from the tangent by formula 9, it remains unclear whether $\lambda$ should be taken in the first or second semicircle. Formula 4 removes the ambiguity by showing that $\lambda$ must be taken in the same semicircle as $u$, since $\cos\varphi$ and $\cos h$ are clearly positive by their nature.”

\par\indent Of course, this addition was needed only for analytical completeness, since there could hardly be any doubt about which value of $\lambda$ is correct. The second value would put the clock error 12 hours off and give the altitude of the pole with the wrong sign. Your extract will also make it easy for anyone to fill in the geometrical discussion given in the programme itself and show what the symbols used mean on the sphere. Professor \textsc{Harding} has also arranged for a German translation of it for the astronomical yearbook.
